\newcommand*{\NAME}{Pieter Nauwelaerts}
\newcommand*{\TITLE}{Extron Version Control} 
\newcommand*{\OFFICE}{UCSC Media Systems Engineering}

\documentclass{article}
\usepackage{graphicx}
\usepackage{hyperref}
\usepackage{lastpage}
\usepackage{fancyhdr}
\usepackage{geometry}
\geometry{margin=1in}
\usepackage{underscore}
\usepackage{subcaption}
\usepackage{fancyvrb}
\usepackage{listings}
\usepackage{caption}
\usepackage{xurl}
\hypersetup{ colorlinks=true,
    linkcolor=blue,
    urlcolor=cyan}
\lstset{ 
    basicstyle=\small\ttfamily,
    columns=flexible,
    breaklines=true,
    numbers=left,
    commentstyle=\itshape\color{gray},
    keywordstyle=\bfseries\color{keywordsColor},
    language=Python,
    showspaces=false,
    showstringspaces=false,
    showtabs=false,
    stringstyle=\color{green!40!black},
    frame=single
}

\usepackage{xcolor}
\definecolor{commentsColor}{rgb}{0.497495, 0.495464, 0.497464}
\definecolor{keywordsColor}{rgb}{0.500000, 0.000000, 0.635294}
\definecolor{stringColor}{rgb}{0.558215, 0.000000, 0.135316}

\title{\TITLE}
\author{\NAME}
\date{\OFFICE}

\begin{document}
\pagestyle{fancy}
\fancyfoot{}
\fancyhead{}
\fancyfoot[L]{\TITLE\ -- \OFFICE}
\fancyfoot[R]{\thepage}

\maketitle

\section*{Introduction}\label{sec:Intro}

This is the collection point of every change made to the Extron ControlScript Code. It is a detailed changelog containing 
version numbers, information on what changed, and background information on why changes needed to be made. It collects 
all of the Git commit messages into one comprehnsible document with more details. For more in depth line by line changes, go 
to the Git repository here: \url{https://git.ucsc.edu/pnauwela/extron}. Each commit has a line by line comparison of what 
was added or removed. The structure of this document has the most recent commits listed at the very top. Version number structure 
is as follows: \lstinline|vX.Y.Z|. X refers to the system components like Switchers or Projectors. When these items are changed out, 
the number will tick up by 1. This number will aslo change for large code base structure changes, including file structure changes, 
or function overhauls. Y refers to big code or UI concepts like adding a new page, or changing code structures and functions. 
Z refers to small line changes, like changing the range on an audio level, or making something work slightly different then it 
did in the previous versions. Another thing to note is that not every room will be on the same version at the same time. Some 
rooms don't get changes in one version but do get the next. Even if a room is at 1.4.3, that doesn't mean it got the change 
from 1.4.2 added to it if 1.4.2 was made for another room. A room will only get a new version number if that version change 
was made to it. Rooms where this happens are noted specifically with their changes in each section. If rooms share a version 
number with different features, it will be marekd. 

\section*{v2.1.0 (Prepared)}

Version 2.1.0 is a complete overhaul of the old code. Each room will get its own version including a few major points:

\begin{itemize}
    \item Password page, and admin passcode function changes.
    \item Inline statements used in button events. 
    \item Button to popup name dictionaries whenever possible.
    \item Simplification of list indexing.
    \item Combining slider files with advanced settings files.
    \item Text files with descriptions of file structure and methods. 
    \item Redundant code removed.
\end{itemize}

The different types of rooms have some variation of implementation, but each room gets these changes. 

\subsubsection*{Passwords}

The password touch panel and control files, as well as the fucntions in the advanced settings page files have been simplified.
Each includes a new function \lstinline|clear_code()| which resets global variables and clears the label. This function is 
used by 3 button events, removing redundant code and decluttering functions. 

Password button definitions are now simplified by changing how the correct input is calculated. Previously, the name of the 
button in the GUI needed to be changed to match its number, making changing a GUI difficult if one needed to be copied over
for a quick deployment. Now, by using the list index for the total button set, with \lstinline|[0 - 9]| in order, the list 
index can be directly called and added to the entered password for comparison. This change also simplifies the code, removing
the need for the \lstinline|button.Name| function call. 

Other simplifications such as changing button states to inline if-statements simplify function design. 

\subsubsection*{Inline Statements}

All touch panel files now feature inline if statements for setting button feedback. 

\subsubsection*{Popup Dictionaries}

Open and close buttons for popup pages are now tied to a dictionary for their page reference. Previously, buttons and page 
names had seperate lists. Now, they have been combined into a dictionary, making opening and closing a page possible through 
one simple command. The exception is the Mac and Windows help pages. Two keys for one item isn't possible, so the open and 
close buttons would need their own index refering to the same popup page. I deemed this to be cluttered as the popup names 
are long. Instead, the open and close buttons are in a list, with \lstinline|[mac, windows]|. The button event is for both
sets, checking if it's in the open or close first, then calling the correct popup function, using the index of the current 
button as the index for the correct popup reference. This is simplified, but is possible to be simplified more with time. 

A notable replacement of lists with dictionaries is in rooms with Biamps. Now, Biamp commands, specifically tags, are sent 
with a dictionary reference. 

\subsubsection*{Simplification of List Indexing}

\subsubsection*{Sliders}

\subsubsection*{Text Files}

\subsubsection*{Redundancy}


\section*{v1.6.4 (Deployed)}

\section*{v1.6.3}

\section*{v1.6.0}

\section*{1.5.6}

\section*{v1.5.5 (Deployed)}

Version 1.5.5 is a version for every room. This version changes all projectors from being connected with 'Power' to 
being connected with 'LampUpdate'. This allows GVE to be updated with the current lamp hours of every room. Power commands 
are sent often enough to be updating the informaton on GVE, and their status is already being sent over. We want GVE to have 
Lamp Hours in the information section, so this is how that will be accomplished. The code is not currently in a room but is ready
to be uploaded. Rooms that are not on larger version 1.5 will be incrememnted in their current version by 1 (for example KRS 
3401 is on v1.4.0, so now is on 1.4.1)

\section*{v1.5.4 (Deployed)}

Version 1.5.4 is a quick fix for KRS3105 and KRS3201. The version was created to fix a strange matrix tie issue. When looking
into the XTP ties, we found the audio from a selected source remained tied to the YuJa output. This issue is not a major problem 
because we don't pull audio from the HDMI inputs to the YuJa, but it still is something that needed to be fixed. The error is caused
by any source being tied to YuJa with Audio/Video. When the shutdown command is sent, the Video portion of that tie gets taken 
over by the Cynap, but the Audio portion remains from whichever source was previously used. When the system is used again, the 
selected source takes over both Audio and Video again. Now, when the system is shutdown, both Audio and Video are tied to the cynap,
removing this issue. 

\section*{v1.5.3}

Version 1.5.3 was for KRS3105 and KRS3201. The version removed YuJa blank image ties from the YuJa output on both shutdown and
startup. In an edge case not originally considered, a user can turn on the video mute before shutting down the system. The 
shutdown and startup were set up to handle this case for the projector outputs, but did not unmute the YuJa outputs. 
This resulted in an error where the Video Mute button was visually turned off and the projector would display content, but the 
YuJa device would not receive any output. Version 1.5.3 fixed this issue for the two rooms. This is not an issue for the smaller 
room type and the IN1808 switchers use a Global Video Mute which gets disabled on startup and shutdown. 

\section*{v1.5.2}

Version 1.5.2 was for KRS3105 and KRS3201. The main change this version made was to add in new matrix tie commands to mute 
YuJa outputs when a video mute button was pressed. The button now mutes the corresponding projector and YuJa output while 
leaving the confidence monitor unmuted. This allows insrtuctors to type in passwords or sensitive information without
students seeing in the room on on recordings. This version was required because the original video mute changes made during 
the rework of Shutdown and Startup did not include YuJa mutes. 

\section*{v1.5.1}

This version is the deployment version of v1.5.0 in which a new login popup was added while waiting for the system to initialize.
A small fix was needed for deployment to the room.  

\section*{v1.5.0}

Version 1.5.0 was developed because of a time delay
The Startup function includes a 1 second delay because of the number commands it contains. This meant the login page was held 
for a second but the passcode cleared making it look like an incorrect pin was put in. This version adds a new UI element which 
lets the user know the system is starting up. This prevents a user from trying to enter the code in the delay. This version is 
a UI change meaning any room needing this version needs a full deployment. A code-only deployment will still work for any future
versions, if the UI change hasn't been made yet. The code to show the popup is included in every room past v1.5.0 even if they 
haven't gotten the full deployment. 

\section*{v1.4.1} 

Version 1.4.1 is a fixed version of v1.4.0. The version added a few new projector commands to the startup and shutdown 
routines created by v1.4.0. The version also removed some of the default audio level variables, instead using direct 
numbers and not allowing them to be changed dynamically. On startup and shutdown, the system sets both mic and program
to an exact default which is fairly low in the range. Users can still raise the audio on the main page. v1.4.1 was initially
prepped for THI001.

\section*{v1.4.0}

\section*{v1.3.2}

\section*{v1.3.0}

\section*{v1.2.2}

\section*{v1.2.1}

\section*{v1.2.0}

\section*{v1.1.1}

\section*{v1.1.0}

\section*{v1.0.1}

\section*{v1.0.0}


\end{document}
